\section{One Taylor center and one sign suffice}\label{sec:cayley}
An ordinary Taylor disk centered at $5$ misses points just to the right of
the critical line at large height. A rational change of variable sends
\emph{the entire right half-plane} to the unit disk. At every positive
real radius the transformed function can already be evaluated by an
absolutely convergent series. This extra information lets us discard one
sign of every coefficient. We give closed integer tests, including a version
requiring only one precision per coefficient.

\subsection*{The data and the six tests per family}
Let $F$ be the coded monic polynomial of degree $d\ge1$, with the height
parameter $H$ of the input packet, and put $C=1000(d+1)^4(H+2)^2$.
For reference, write
\[
 h_F(p,k)=\deg\gcd(\bar F,X^{p^k}-X),\quad
 B(y)=\left\lfloor\log_2\prod_{p^k\le y}p^{h_F(p,k)}\right\rfloor,
 \quad a_n=B(n)-B(n-1).
\]
Primes $p$ and positive integers $k$ occur in the product, with $B(0)=0$.
Define $b_n=a_n-2a_{n/2}$ at even $n$, and $b_n=a_n$ at odd $n$.
For a cutoff $N\ge1$ let
\[
 J_n=\sum_{i=1}^{N(n-1)}\left\lfloor\frac{N^2}{N+i}\right\rfloor,
 \qquad A_k=\sum_{n=1}^N b_n J_n^k (N!/n)^5,
 \qquad D_k=k!N^{2k}(N!)^5 .
\]
All denominators below are positive. At each order $m\ge1$ put
\begin{align*}
 T_{m,N}
   &=\sum_{k=1}^m (-9)^k\binom{m-1}{k-1}
       \frac{m!}{k!}N^{2(m-k)} A_k,\\
 E_m&=81(d+1)10^{m-1},\qquad V_m=13500C(m+1)^4 .
\end{align*}
Choose \emph{one} of the upper, lower and two-sided comparisons
\begin{equation}\label{eq:cayleytest}
  NT_{m,N}\le D_m(N V_m+E_m),\qquad
  -NT_{m,N}\le D_m(N V_m+E_m),\qquad
  N|T_{m,N}|\le D_m(N V_m+E_m).
\end{equation}
Quantify over every $m\ge1$. Independently choose between testing
\emph{every} $N\ge1$ and testing only $N=E_m$.
For RH use $F=X-1$, for GRH use every $F=X^q-1$ with
$q\ge1$, and for ERH use every polynomial packet. Thus no zero, contours,
real constants, choice functions, or numerical witnesses occur in a test.
A signed inequality is compiled by moving its two nonnegative sums to
opposite sides and testing whether their truncated difference is zero.

\paragraph{Theorem.}
Each of these six universal sentences, interpreted in the standard
integers, is equivalent to the corresponding analytic hypothesis. A single
sign at the prescribed precision is sufficient. In the two-sided direction
the coefficient estimate loses at most one power of $m+1$ when recovered
from one sign. The explicit bound is given below.

\subsection*{From off-line zeros to poles of one germ}
Use $H_F(s)=\sum_{n\ge1} b_n n^{-s}$ for $\Re s>1$.
We recall the analytic inputs from the attached prime-product construction.
If $J_F(s)=\sum_{p,k}h_F(p,k)(\log p)p^{-ks}$ and
$\Psi_F(x)=\sum_{p^k\le x}h_F(p,k)\log p$, Abel summation gives
\[
 H_F(s)=(1-2^{1-s})\left(\frac{J_F(s)}{\log 2}
 +s\int_1^\infty
 \left(B(\lfloor x\rfloor)-\frac{\Psi_F(x)}{\log 2}\right)x^{-s-1}\,dx\right).
\]
The discrepancy inside the integral is bounded by $1$, so the integral
is holomorphic on $\Re s>0$. Outside finitely many bad primes, the local
factors identify $J_F$ with minus the sum of logarithmic derivatives of
the Dedekind zeta functions belonging to the distinct irreducible factors
of $F$ over $\mathbb Q$. At a bad prime the correction is holomorphic on
$\Re s>0$ by its geometric-series expression. At a common zero the
residues of the logarithmic derivatives add with the same sign. In
$1/2<\Re s<1$, the factor $1-2^{1-s}$ is nonzero; it cancels the
possible pole at $s=1$. Consequently an off-line zero on the right gives
a genuine pole. The functional equations reflect zeros on the left.
The cyclotomic factorization accounts for all primitive Dirichlet
characters, and an integral primitive element gives a packet for every
number field. These are classical identification steps.

The other input is the numerical bound under the corresponding hypothesis,
\begin{equation}\label{eq:cayleysum}
   \left|\sum_{n\le x} b_n\right|\le 100 C\sqrt{x}(1+\log x)^2
   \qquad(x\ge1).
\end{equation}
Indeed the prefix is $B(\lfloor x\rfloor)-2B(\lfloor x/2\rfloor)$.
At an even cutoff $2n\ge200$ this is the original dyadic gap with bound
$C\sqrt n(1+\lfloor\log_2(2n)\rfloor)^2$. An adjacent odd cutoff
adds at most $d\log_2(2n+1)+1$. The elementary bound
$B(y)\le d y\log_2 y$ covers $1\le y<200$. These give
\eqref{eq:cayleysum} with the displayed constant.

Use the change of coordinate
\[
 s(z)=\frac{5+4z}{1-z},\quad z(s)=\frac{s-5}{s+4},\quad
 \Re s(z)-\tfrac12=\frac{9(1-|z|^2)}{2|1-z|^2}.
\]
In particular $s(0)=5$ and the map is biholomorphic between the unit disk
and $\Re s>1/2$. Write $f(z)=H_F(s(z))=\sum_{m\ge0}v_m z^m$ locally
at zero. These $v_m$ are real. Introduce the ordinary Taylor coefficients
$u_k=\sum_{n\ge1} b_n(\log n)^k/(k!n^5)$ of $H_F(5-w)$. Near zero,
$5-s(z)=-9z/(1-z)$. Expanding $(1-z)^{-k}$ gives
\[
 v_0=u_0,\qquad
 v_m=\sum_{k=1}^m(-9)^k\binom{m-1}{k-1}u_k .
\]
For example, $v_1=-9u_1$, $v_2=-9u_1+81u_2$ and
$v_3=-9u_1+162u_2-729u_3$.

\paragraph{Certified finite approximation.}
The bound $|u_k-A_k/D_k|\le9(d+1)/N$ is uniform in $k$.
For completeness, comparing the right rectangle sum for $1/x$ on
$[1,n]$ with mesh $1/N$ to its integral, then allowing $1/N^2$ per
floor, gives $0\le\log n-J_n/N^2\le n/N$.
For $k\ge1$ the derivative of $t^k/k!$ there is bounded by $n$.
Also $|b_n|\le3(d+1)n$ and $(\log n)^k/k!\le n$. The first $N$
terms therefore cost at most $3(d+1)\sum n^{-2}/N\le6(d+1)/N$,
and the tail costs at most $3(d+1)\sum_{n>N}n^{-3}\le3(d+1)/N$.
For $k=0$ the prefix has zero error. The binomial theorem gives
$\sum_{k=1}^m9^k\binom{m-1}{k-1}=9\cdot10^{m-1}$. Clearing the
common denominator now gives
\[
              |v_m-T_{m,N}/D_m|\le E_m/N .
\]
At $N=3$ the RH packet gives approximations $-299/1944$, $383/1944$,
$-121/5832$ for the first three coefficients. These are small-cutoff
exact arithmetic examples; the certified errors here are large.

\paragraph{Forward implication.}
Under \eqref{eq:cayleysum}, partial summation gives
$H_F(s)=s\int_1^\infty S(x)x^{-s-1}\,dx$ for $\Re s>1/2$,
where $S(x)=\sum_{n\le x}b_n$. On $|z|=r<1$ we have
$|s|\le9/(1-r)$ and $\delta=\Re s-1/2\ge(9/4)(1-r)$.
Substituting $x=e^t$ gives the integral majorant
$100C|s|(\delta^{-1}+2\delta^{-2}+2\delta^{-3})$.
Thus $|f(z)|\le4500C(1-r)^{-4}$. At $r=m/(m+1)$, Cauchy's
estimate and $(1+1/m)^m<3$ give $|v_m|\le V_m$. Adding the finite
approximation error proves every assertion in \eqref{eq:cayleytest}.

\subsection*{Recovering the missing sign}
A strict one-sided estimate on the growth order is generally insufficient
for an arbitrary power series. Here the positive real interval is special:
$s(x)\ge5$ when $0\le x<1$. We prove the requisite positivity
argument directly and give finite translation bounds.

It is enough to assume the upper comparison at $N=E_m$. It gives
$v_m\le V_m+2$; the error $1$ occurs once in the test and once in passing
to the true coefficient. For the lower comparison replace $f$ by $-f$.
Let $K=13500C+2$, write $g=f$ in the upper case and $g=-f$ in the
lower case, and denote the coefficients of $g$ by $g_j$. Then
$c_j=K(j+1)^4-g_j\ge0$, including $j=0$ since $|u_0|\le6(d+1)$.
Put
\[
 G(z)=K\sum_{j\ge0}(j+1)^4z^j-g(z).
\]
The first term is rational on the disk, obtained by applying
$(1+z\,d/dz)^4$ to $(1-z)^{-1}$. The definition of $g$ by its
Dirichlet series supplies analytic neighborhoods along $[0,1)$.

We claim $c_m r^m\le G(r)$ for rational $0<r<1$. Subtract the finite
prefix $P(z)=\sum_{j=0}^m c_j z^j$. All Taylor coefficients of $G-P$
at $0$ are nonnegative. At a real center $x$, translating its Taylor
coefficients through a positive distance $h$ inside a disk of convergence
uses only nonnegative weights $\binom{j+k}{j}h^k$. Thus nonnegativity
passes along finitely many overlapping disks to $r$, proving
$G(r)-P(r)\ge0$. Since $P(r)\ge c_m r^m$, this is the claim.
This also explains the classical Pringsheim argument in this setting.

For an explicit mesh take $R=(1-r)/8$ and positive steps $h\le R/4$.
On $|z-x|\le R$, $0\le x\le r$, we have $\Re s(z)\ge4$. Indeed
the real part of $1/(1-z)$ is at least $1/(1-x+R)$ on this disk, and
$s=9/(1-z)-4$. Hence $|g(z)|\le3(d+1)\sum n^{-3}\le6(d+1)$
there. Since $(j+1)^4\le24\binom{j+4}{4}$,
\[
  G(r)\le 24K(1-r)^{-5}+6(d+1).
\]
At $r=m/(m+1)$ we obtain $c_m\le 3[24K(m+1)^5+6(d+1)]$ and hence
$|v_m|\le100K(m+1)^5$. The Taylor series now converges absolutely on
the whole disk. It agrees near zero with the meromorphic continuation
of $H_F(s(z))$ and by uniqueness cannot hide any pole there. The pole
comparison proves the converse for each family.

\paragraph{How many derivatives and bits does positivity require?}
This detail matters if induction is to be restricted to bounded,
decidable computations. All derivatives at a point need never be
propagated in a single induction invariant. Let $L(z)=G(z)-P(z)$.
A uniform integer bound $M\ge1$ on $|L|$ on the mesh disks is obtained
by rounding up
\[
 1+24K\bigl(1-r-R\bigr)^{-5}+6(d+1)
       +\sum_{j=0}^m \bigl(K(j+1)^4+54(d+1)10^j\bigr).
\]
Indeed $|u_k|\le6(d+1)$ follows by absolute summation, hence the
unconditional rough bound on $|v_j|$ used in this expression.
Let $w_j(x)=R^j L^{(j)}(x)/(j!M)$. Cauchy's estimate gives
$|w_j(x)|\le1$. Translation from $x$ to $x+h$ yields
\[
 w_j(x+h)=\sum_{k\ge0}\binom{j+k}{j}(h/R)^k w_{j+k}(x).
\]
The total weight is at most $2^{j+1}$ and the tail after $k=b$ is at most
$2^{j-b}$, using $\binom{j+k}{j}\le2^{j+k}$ and $h/R\le1/4$.
To prove $w_j(x+h)\ge-2^{-e}$ for every $0\le j\le J$, retain
$0\le k\le J+e+1$. It is enough to know
$w_i(x)\ge-2^{-(e+J+2)}$ for $i\le2J+e+1$. The weighted error
and the tail each cost at most $2^{-e-1}$. Start at the endpoint with
$(J,e)=(0,t)$ and work backwards by
\[
             (J,e)\longmapsto(2J+e+1,e+J+2).
\]
There are at most $\lceil32r/(1-r)\rceil$ steps. Every requested
entry at zero is nonnegative, so $L(r)\ge-M2^{-t}$ for arbitrary
$t$. Taking this as the definition of closed order on rational Cauchy
names gives $L(r)\ge0$. The mesh, order, precision, and sums are
primitive recursive in the rational data. This argument isolates the
finite propagation from the analytic certificates that justify Taylor
translation and its tails.

\paragraph{Formal scope.}
\texttt{CayleyPackets.lean} constructs the 18 closed object formulas.
\texttt{matrix\_reading} and \texttt{statement\_reading} identify their
integer meanings. \texttt{weightSum} proves the exact binomial amplification.
\texttt{translation\_step} proves the finite weighted error inequality over
rationals and \texttt{finite\_propagation} verifies the backwards budget
recursion. The object derivations \texttt{scheduled} and
\texttt{certified\_failure} identify, respectively, equivalent single-index
cube tests and least-radius positive counterexample certificates in iPRA.
The holomorphic identification, Cauchy estimates, and analytic re-expansion
on the mesh are human proofs. The Lean development supplies neither an
internal iPRA derivation from the prime-gap sentence to the one-sided tests
nor an internal derivation of the analytic limit identities.
